\section{Introduction}

AI agents are rapidly becoming more capable. They can now solve tasks as complex
as resolving real-world GitHub issues \cite{yang2024sweagent} and real-world
email organization tasks \cite{roth2024google}. However, as their capabilities
for benign applications improve, so does their potential in dual-use settings.

Of the dual-use applications, hacking is one of the largest concerns
\cite{lohn2022will}. As such, recent work has explored the ability of AI agents
to exploit cybersecurity vulnerabilities \cite{fang2024llm, fang2024llm2}. This
work has shown that simple AI agents can autonomously hack mock
``capture-the-flag'' style websites and can hack real-world vulnerabilities when
given the vulnerability description. However, they largely fail when the
vulnerability description is excluded, which is the \emph{zero-day exploit}
setting \cite{fang2024llm2}. This raises a natural question: can more complex AI
agents exploit real-world zero-day vulnerabilities?

In this work, we answer this question in the affirmative, showing that
\emph{teams} of AI agents can exploit real-world zero-day vulnerabilities. To
show this, we develop a novel multi-agent framework for cybersecurity exploits,
extending prior work in the multi-agent setting \cite{liu2023dynamic,
chen2023autoagents, zhang2023building}. We call our technique \sn, which (to our
knowledge) is the first multi-agent system to successfully accomplish meaningful
cybersecurity exploits.

Prior work uses a single AI agent that explores the computer system (i.e.,
website), plans the attack, and carries out the attack. Because all highly
capable AI agents in the cybersecurity setting at the time of writing are based
on large language models (LLMs), the joint exploration, planning, execution is
challenging for the limited context lengths these agents have.

We design \emph{task-specific, expert} agents to resolve this issue. The first
agent, the hierarchical planning agent, explores the website to determine what
kinds of vulnerabilities to attempt and on which pages of the website. After
determining a plan, the planning agent dispatches to a team manager agent that
determines which task-specific agents to dispatch to. These task-specific agents
then attempt to exploit specific forms of vulnerabilities.

To test \sn, we develop a new benchmark of recent real-world vulnerabilities
that are past the stated knowledge cutoff date of the LLM we test, GPT-4. To
construct our benchmark, we follow prior work and search for vulnerabilities in
open-source software that are reproducible. These vulnerabilities range in type
and severity.

On our benchmark, \sn achieves a pass at 5 of 42\%, within 1.8$\times$ of a GPT-4
agent with knowledge of the vulnerability. Furthermore, it outperforms
open-source vulnerability scanners (which achieve 0\% on our benchmark) and a
single GPT-4 agent with no description. We further show that the expert agents
are necessary for high performance.

In the remainder of the manuscript, we provide background on cybersecurity and
AI agents (Section~\ref{sec:background}), describe the \sn
(Section~\ref{sec:arch}), our benchmark of real-world vulnerabilities
(Section~\ref{sec:benchmark}), our evaluation of \sn (Section~\ref{sec:eval}),
provide case studies (Section~\ref{sec:case-studies}) and a cost analysis
(Section~\ref{sec:cost}), describe the related work (Section~\ref{sec:rel-work})
and conclude (Section~\ref{sec:conclusion}).
